\begin{frame}
\frametitle{Breathing}
%\begin{columns}[c] % The "c" option specifies centered vertical alignment while the "t" option is used for top vertical alignment

%\column{.3\textwidth} % Left column and width
%\includegraphics[width=\linewidth]{Sitting_chair_side.png}
%\column{.7\textwidth} % Right column and width
\begin{itemize}
\item[-] The \structure{breath} is of utmost importance for our \structure{health and well--being}.
\item[-] The nice things about is, that we can have an \structure{influence} on our \structure{breath}, therefore on \structure{our health and well--being}.
\item[-] When we are \structure{focused, in joy, bored, in a hurry, physically active, angry, getting agitated or sad}, our breath goes along accordingly.
\item[-] Our \structure{breath adapts} to our state of being.
\item[-] \structure{Each unconscious function of our body} is in turn connected to the breath.
  \item[-] Breathing, depending on how we train it, \structure{builds up strength} and brings  missing or excessive things into \structure{equilibrium}.
  \end{itemize}
% Write on
%\end{columns}
\end{frame}
%--------------------------------------------------------------------------------------------------------------

\begin{frame}
\frametitle{Perception of the Body}
%\begin{columns}[c] % The "c" option specifies centered vertical alignment while the "t" option is used for top vertical alignment

%\column{.3\textwidth} % Left column and width
%\includegraphics[width=\linewidth]{Sitting_chair_side.png}
%\column{.7\textwidth} % Right column and width
\begin{itemize}
\item[-] Conscious breathing also helps to better \structure{perceive your body} and therefore \structure{yourself}.
\item[-] Breathing is the \structure{instrument}, that you can lead through your body in a \structure{targeted and conscious way}.
\item[-] Your breath also leads you into a \structure{world of love}, as you are \structure{connected through your breath} with everybody and everything.
\item[-] In the framework of stress regulation, the breath is \structure{most immediately active tool}.
  \item[-] You can \structure{use it everywhere} and it's a natural tool --- and it also doesn't cost anything.
    \end{itemize}
% Write on
%\end{columns}
\end{frame}
%--------------------------------------------------------------------------------------------------------------

\begin{frame}
\frametitle{Basic principles}
%\begin{columns}[c] % The "c" option specifies centered vertical alignment while the "t" option is used for top vertical alignment

%\column{.3\textwidth} % Left column and width
%\includegraphics[width=\linewidth]{Sitting_chair_side.png}
%\column{.7\textwidth} % Right column and width
\begin{description}
\item[The inhale] \structure{builds up}, gives \structure{strength} and puts you \structure{back on your feet}.
\item[The exhale] \structure{relaxes}, calms down and \structure{loosens tension}.
  \end{description}

\begin{itemize}
\item[-] If you want to \structure{refresh}, build yourself up and get \structure{strength} --- \structure{prolong the inhale} and the \structure{pause} after the inhale.
\item[-] If you want to \structure{relax} and come to \structure{rest} --- \structure{slow down} and deepen the \structure{exhale}.
  \item[-] If you want to create an \structure{inner balance} --- pay attention to the inhale and exhale being \structure{regular}.
\end{itemize}

Don't force anything by doing that. \structure{Observe the rhythm} in a relaxed manner. When the \structure{thoughts are wandering} elsewhere, just bring them \structure{back to the breath}.
% Write on
%\end{columns}
\end{frame}
%--------------------------------------------------------------------------------------------------------------


\begin{frame}
  \frametitle{The optimal breathing rhythm}

  According to new scientific findings, an \structure{optimal breathing rhythm} looks like:

  6 breaths a minute, so therefore \structure{10 seconds per breath}, distributed as follows:

  \begin{description}
  \item[1 -- 4] inhale
  \item[5 -- 8] exhale
    \item[9 -- 10] pause on empty lungs
    \end{description}
\end{frame}
%--------------------------------------------------------------------------------------------------------------

\begin{frame}
  \frametitle{Tips for your breathing}

  \textbf{Flat and only superficial breathing (upper chest only)}

    \vspace{5 mm}

    While exhaling, just press with \structure{both hands on the belly} and then let the inhale freely stream into your body.
    You can achieve the same, by pressing the \structure{hands sideways onto your rips (flank)} or on the tips of your lungs (upper lung space).
    
    And you already completed a \structure{perfect full breath cycle}! Isn't that just easy as pie?

    \end{frame}
%--------------------------------------------------------------------------------------------------------------

\begin{frame}
  \frametitle{Tips for your breathing}

  \textbf{Leading the full breath into your body}

  \vspace{5 mm}

    By the way, you can lead your breath \structure{into any place into your body} this way. Try it!
    
    If you have \structure{pain} somewhere in your body, consciously lead your \structure{breath into that part} of your body.
    You will notice how much that reliefs.
    You become \structure{your own healer}.
  
\end{frame}
%--------------------------------------------------------------------------------------------------------------


\begin{frame}
  \frametitle{Tips for your breathing}

  \textbf{Not being able to easily exhale}

    \vspace{5 mm}

  In this case, there's also a simple exercise to help you:
  
  \structure{Exhale} while making a `\structure{`shshsh'' sound}.

    \vspace{5 mm}
  
  Your \structure{exhale will prolong}.

  \end{frame}
%--------------------------------------------------------------------------------------------------------------

\begin{frame}
  \frametitle{Tips for your breathing}

  \textbf{A breath workout to optimize your inhale/exhale}

     \vspace{5 mm}

     If you want to train your breath and optimize your inhale and exhale, you can make a \structure{``breathing regimen''},
     practice it \structure{every day 3 times, 5 minutes each}.
  You will quickly \structure{feel better, energized and more alive} and your \structure{capacity to focus and your mood} will get better.

  \end{frame}
%--------------------------------------------------------------------------------------------------------------

\begin{frame}
  \frametitle{Alternating breathing}
  
  \textbf{1 cycle of alternating breathing}
  \vspace{5mm}

  \begin{itemize}
  \item[-] \structure{Close the right nostril} with the right index finger and \structure{inhale} through your \structure{left side}.
  \item[-] \structure{Hold your breath} for a short while.
  \item[-] Now, \structure{close the left nostril} with the left index and \structure{exhale} through your \structure{right side}.
  \item[-] \structure{Inhale} on your \structure{right side}.
  \item[-] \structure{Hold your breath} for a short while.
  \item[-] Now, \structure{close the right nostril} and \structure{exhale} through your \structure{left side}. And so on...
\end{itemize}
  
  
  \end{frame}
%--------------------------------------------------------------------------------------------------------------
\begin{frame}
  \frametitle{Alternating breathing}

  \textbf{Alternating breathing --- Execution and Benefits}
  \vspace{5mm}

  
  \begin{itemize}
    
  \item[-] \structure{Continue} this alternating breath, \structure{6 -- 12 times}.
    \item[-] To \structure{finish}, take at least \structure{6 deep breaths},  through \structure{both nostrils and consciously perceived}.
\end{itemize}

  \vspace{5 mm}
  
  The practice of alternating breath helps you to get \structure{back into equilibrium},
  homeostasis, achieve a \structure{stable mood} and give you \structure{strong nerves}, to be ready for whatever comes your way.
  
  \end{frame}
%--------------------------------------------------------------------------------------------------------------

\begin{frame}
  \frametitle{Breathing to regulate energy level}

  \textbf{Breathing for tiredness and bad mood}
  \vspace{5mm}

  
  \begin{itemize}
    
  \item[-] \structure{6 --12 times inhale right} and \structure{exhale left}.
    \item[-] Prolong the \structure{pause after each inhale}.

  \end{itemize}

  \vspace{5 mm}
  
  \textbf{Breathing for restlessness and nervousness/anxiety}
  \vspace{5mm}

  
  \begin{itemize}
    
  \item[-] \structure{6 --12 times inhale left} and \structure{exhale right}.
    \item[-] Prolong the \structure{pause after each exhale}.

    \end{itemize}
  
  
  \end{frame}
%--------------------------------------------------------------------------------------------------------------

%--------------------------------------------------------------------------------------------------------------

\begin{frame}
  \frametitle{Effects of breathing exercises}

  These breathing exercises \structure{act immediately} on your body.

  The effects hold for a longer time and get more pronounced, if you prescribe yourself a \structure{regimen}.

  \vspace{3 mm}
  
  That means, you practice them \structure{4 times daily for 5 minutes each}, over a period of about \structure{3 weeks}.

  \vspace{10 mm}

  \textbf{Source:} translated from: ``Yoga ganz einfach'' by Gertrud Hirschi (\url{https://gertrudhirschi.ch}).

  
  
  \end{frame}
%--------------------------------------------------------------------------------------------------------------
